\documentclass[12pt]{article}
\usepackage{amsmath,amssymb}

\begin{document}
\section*{{2025 USAMO Problems/Problem 4}}
Source: AoPS Wiki

\subsection*{Solution}
Let AP intersects BC at D. Extend FC to the point E on the circumcircle $\omega$ of $FAP$ . Since $H$ is the orthocenter of $\Delta ABC$ , we know that $HD = DP$ or $HP = 2HD$ , and $AH \cdot HD = CH \cdot HF$ . Next we use the power of H in $\omega$ : $AH \cdot HP = CH \cdot HE$ . These relations imply that $HE = 2HF$ . Hence $C, D$ are midpoints of $HE, HP$ respectively. By midline theorem, $CD // EP$ . Since $AD \perp CD$ , we have $AD \perp EP$ . This implies that $\angle APE = 90^{\circ}$ . Consequently, $AE$ is the diameter of $\omega$ . Let $G$ be the midpoint of $AE$ which is also the center of $\omega$ . $G,C$ are midpoints of $AE, EH$ respectively. By the midline theorem again, we have $GC//AH$ , consequently, $GC \perp BC$ . This implies that $GC$ is the perpendicular bisector of the chord $XY$ hence $C$ is the midpoint of $XY$ . ~ Dr. Shi davincimath.com

\end{document}
